% CSE 540 (Fall 2026, Group 18) — Final technical report.
% Compile on Overleaf: this file is self-contained IEEEtran conference format;
% no external figures or .bib files are required.
%
% Team review notes live in docs/final-report/README.md.

\documentclass[conference]{IEEEtran}
\usepackage{cite}
\usepackage{booktabs}
\usepackage{array}

\begin{document}

\title{Chain-Anchored Provenance and Compliance for Electronics Hardware Bills of Materials}

\author{\IEEEauthorblockN{Pranav Mandava}
\IEEEauthorblockA{\textit{Group 18} \\ \textit{CSE 540: Engineering Blockchain Applications} \\ pranav.techiegeek@gmail.com}
\and
\IEEEauthorblockN{Suhas Gajula}
\IEEEauthorblockA{\textit{Group 18} \\ \textit{CSE 540: Engineering Blockchain Applications} \\ suhasgajula248@gmail.com}
\and
\IEEEauthorblockN{Thushar Reddy Vemula}
\IEEEauthorblockA{\textit{Group 18} \\ \textit{CSE 540: Engineering Blockchain Applications} \\ vemulathusharreddy@gmail.com}
\and
\IEEEauthorblockN{Varshini Akula}
\IEEEauthorblockA{\textit{Group 18} \\ \textit{CSE 540: Engineering Blockchain Applications} \\ akulavarshini@gmail.com}
\and
\IEEEauthorblockN{Venkata Pradeep Reddy Venuthurla}
\IEEEauthorblockA{\textit{Group 18} \\ \textit{CSE 540: Engineering Blockchain Applications} \\ pradeepreddyvv@gmail.com}
}

\maketitle

\begin{abstract}
Electronics hardware bills of materials (HBOMs) span companies that never share records: tier-2 vendors hold component certificates, tier-1 assemblers hold board definitions, and distributors hold custody logs. A buyer assembling a workstation cannot trace which batches entered a board, which compliance documents cover them, or who handled the device last. We present a Solidity prototype that makes this fragmented record a single shared one. Component batches, compliance evidence, board membership, and a two-step custody handover live in one role-gated contract; documents themselves stay off-chain on IPFS, anchored by a keccak256 hash stored at registration time and re-verifiable by anyone. Auditors attest hashes under an on-chain separation-of-duties rule that forbids self-attestation, and a single read returns the exact missing documents for any board. We measure the full gas profile of the system on a deterministic local chain, price three storage-design alternatives (word identifiers, a Merkle-root HBOM, flat role checks), and use those numbers to justify the design: per-certificate records dominate cost (69\% of a 28.9M-gas seeded world) but are what makes exact gap enumeration possible. We close with what the system deliberately cannot prove---physical authenticity---and the privacy, interoperability, and governance constraints a real deployment would face.
\end{abstract}

\begin{IEEEkeywords}
blockchain, provenance, supply chain, smart contracts, HBOM, compliance, RoHS, REACH, IPFS
\end{IEEEkeywords}

\section{Introduction}
An assembled electronics board is a stack of components from companies that do not talk to each other. Each tier keeps its own records: the tier-2 supplier knows which reel of SSDs shipped with which REACH declaration; the tier-1 assembler knows which batches entered which board serial; the distributor knows who held the box last. No single one of them, and certainly not the end buyer, can answer the question that matters: \emph{for this board, which documented, attested evidence covers every component, and who has handled it?}

The cost of not answering it is well documented. Counterfeit integrated circuits evade detection precisely because the paper trail is fragmented across tiers~\cite{guin}, and compliance regimes---RoHS, REACH, FCC---are enforced through documents that a buyer must request, trust, and file. Prior work applies blockchain provenance to electronics supply chains~\cite{cui} and product-ownership management~\cite{toyoda}; both conclude that the binding constraint is not the ledger but the linkage between physical items, documents, and custody events.

This paper describes the design, implementation, and measured cost of \texttt{HBOMProvenance}, a course-scale prototype that makes that linkage the primary artifact:

\begin{itemize}
\item \textbf{One contract, five stakeholders.} Suppliers register batches and attach document commitments; assemblers link batches into boards; auditors attest or revoke documents; custodians move boards through a two-step handover; buyers read everything without a wallet.
\item \textbf{Hash-first evidence.} Files are hashed in the browser before anything is signed; the contract stores the hash, type, and IPFS content identifier (CID). An auditor's attestation therefore binds a \emph{specific document content}, not a claim.
\item \textbf{Exact gaps, on-chain.} Compliance verification does not return a boolean; it returns the precise missing \texttt{component:certType} pairs, and the buyer-facing badge renders that answer verbatim.
\end{itemize}

Section~\ref{sec:design} covers the system design, Section~\ref{sec:contract} the contract logic, Section~\ref{sec:impl} the implementation and test suite, Section~\ref{sec:gas} the gas and scalability measurements, and Section~\ref{sec:discussion} the discussion a production system would have to survive.

\section{System Design}\label{sec:design}

\subsection{Stakeholders and Access Control}
The contract separates five participant types through four operational roles plus a role administrator (Table~\ref{tab:roles}). Roles are OpenZeppelin \texttt{AccessControl} roles~\cite{oz}; every grant and revocation is itself a transaction with an event, so governance history is on-chain by construction rather than by convention. Role assignment is deliberately outside the constructor: the deploy script grants roles through visible transactions, and the demo shows each participant receiving theirs.

\begin{table}[!t]
\caption{Roles and their write permissions}
\label{tab:roles}
\centering
\begin{tabular}{@{}l l l@{}}
\toprule
Role & Held by & Permitted writes \\
\midrule
\texttt{SUPPLIER} & tier-2 vendors & register batches; attach documents to own batches \\
\texttt{ASSEMBLER} & tier-1 manufacturers & create board HBOMs from registered batches \\
\texttt{AUDITOR} & regulators, labs & attest or revoke any document \\
\texttt{CUSTODIAN} & logistics, distributors & initiate/accept/cancel custody \\
\texttt{ADMIN} & consortium operator & grant and revoke all roles \\
\bottomrule
\end{tabular}
\end{table}

\subsection{Data Model}
Three structs carry the state. A \texttt{Component} is a registered batch: human-readable string id (e.g.\ \texttt{SSD-SAMSUNG-990PRO-2T-L2409}), datasheet hash and CID, registering supplier, timestamp. A \texttt{Certificate} is one compliance document against one batch: certificate type (RoHS, REACH, FCC), document hash and CID, issuing supplier, attesting auditor, attestation timestamp, and two flags---\texttt{attested} and \texttt{revoked}. A \texttt{Board} is an assembled unit: its member component ids, assembler, current custodian, a pending custodian, and a custody state (\texttt{InCustody} / \texttt{InTransit}).

Nothing is ever deleted. Revocation is a flag flip, not a removal; the certificate history of every batch remains reconstructable from storage for the life of the chain. At course scale this is a virtue rather than a cost---the audit trail is the product.

\subsection{The Product Journey}
The journey the assignment describes maps to eight write functions:

\begin{enumerate}
\item \emph{Creation}: \texttt{registerComponent(id, datasheetHash, cid)} --- a supplier records a batch.
\item \emph{Evidence}: \texttt{addCertificate(id, type, docHash, cid)} --- the same supplier attaches a compliance document.
\item \emph{Validation}: \texttt{attestCertificate(certId)} --- an auditor independently confirms the document hash; \texttt{revokeCertificate} reverses it with a reason string in the event.
\item \emph{Assembly}: \texttt{assembleHBOM(boardId, componentIds)} --- an assembler links registered batches into a board (bounded by \texttt{MAX\_COMPONENTS = 16}).
\item \emph{Shipment and delivery}: \texttt{initiateTransfer} then \texttt{acceptTransfer} (or \texttt{cancelTransfer}) --- a two-step custody handover.
\item \emph{Verification}: \texttt{verifyCompliance(boardId)} --- anyone, without a wallet, reads compliance and the exact gaps.
\end{enumerate}

\subsection{Architecture}
The client is a TanStack Start application deployed to Cloudflare Workers. The browser talks to the chain through a same-origin JSON-RPC proxy (so neither a local node URL nor any RPC key ever reaches the page), and uploads documents through a same-origin IPFS pinning route backed by Pinata---the JWT stays a server-side secret. Files are hashed \emph{before} upload; if pinning is unavailable, the flow degrades to hash-only with a manually typed CID, and nothing about the trust model is lost, since the hash is the anchor.

Two chain environments back the same contract and client. The graded demonstration runs on a local Hardhat node with five deterministic funded accounts, one per role---the entire demo world seeds in one command (~29M gas, offline, repeatable). The same contract is also deployed publicly: \texttt{0x803e3778dd01ae83e9a25cb11942ef7b6063422c} on Base Sepolia (chain 84532), fronted by the Worker at \texttt{hbom-provenance.pranav-techiegeek.workers.dev}, so anyone can verify the same flows through a live URL with two throwaway accounts; separation of duties is preserved there because the auditor key holds only \texttt{AUDITOR\_ROLE}.

\section{Smart Contract Logic}\label{sec:contract}

\subsection{Errors, Not Strings}
Every failure path is a custom error with its offending values: \texttt{AlreadyExists(id)}, \texttt{NotFound(id)}, \texttt{EmptyId()}, \texttt{InvalidHash()}, \texttt{NotComponentSupplier(id, sender)}, \texttt{AlreadyAttested}, \texttt{AlreadyRevoked}, \texttt{SelfAttestation()}, \texttt{WrongCustodyState(expected, actual)}, \texttt{NotCustodian(id, sender)}, \texttt{SelfTransfer()}, \texttt{InvalidAddress()}, \texttt{NoComponents()}, \texttt{TooManyComponents(n, max)}, and \texttt{DuplicateComponentInBatch(id)}. The client decodes the error name and arguments---a wrong-role attempt on the register screen surfaces as \texttt{AccessControlUnauthorizedAccount(0x..., SUPPLIER\_ROLE)} to the user, which doubles as the role-gating demonstration in the demo.

\subsection{Separation of Duties}
Two rules are enforced in \texttt{attestCertificate}: an already-attested certificate cannot be attested again (\texttt{AlreadyAttested}), and the auditor cannot have issued the document themselves (\texttt{SelfAttestation}). Together they guarantee that every accepted document has been confirmed by \emph{two} distinct keys. Revocation is deliberately allowed even before attestation: an auditor inspecting a fraudulent document can reject it outright, and the record of the rejection stays.

Only the supplier that registered a batch may attach evidence to it (\texttt{NotComponentSupplier}), so a competitor cannot poison another vendor's record even while holding the supplier role.

\subsection{Two-Step Custody}
\texttt{initiateTransfer} moves a board to \texttt{InTransit} and names a pending custodian; only that account can \texttt{acceptTransfer}. Two guards make the state machine safe under typos: the recipient must already hold \texttt{CUSTODIAN\_ROLE} (otherwise a mistyped address would strand the board in \texttt{InTransit} forever), and the current custodian cannot transfer to themselves. The original custodian retains an escape hatch: \texttt{cancelTransfer} returns the board to \texttt{InCustody} until the moment the recipient accepts.

\subsection{Compliance Semantics}
\texttt{verifyCompliance} walks every component of a board and requires, for each of the three certificate types, at least one certificate that is attested \emph{and} not revoked. The return value is the list of gaps as \texttt{componentId:certType} strings; the buyer page renders it verbatim. This is the load-bearing design decision of the system: a boolean would be cheaper to consume but would force the buyer back into asking each tier \emph{which} document is missing---the exact interaction this project removes. Section~\ref{sec:gas} prices this decision and shows it is nearly free (the gap-reporting path costs $+0.9\%$ over the clean path).

\subsection{Events and Stateless Clients}
Every state change emits an event carrying the human-readable ids as unindexed strings. The web client rebuilds its entire view by replaying events from block zero, then hydrating records with storage reads---no indexer, no database, no backend state. At course scale (30 batches, 90 certificates) this is a handful of RPC calls; at production scale it is the first component to replace (Section~\ref{sec:future}).

\section{Implementation and Testing}\label{sec:impl}
The contract is Solidity 0.8.28 with the optimizer at 200 runs, inheriting OpenZeppelin \texttt{AccessControl}. Tooling is Hardhat 3 with viem shared between the test suite and the web client. The test suite covers the protocol rules promised in the proposal: the happy path; unauthorized writes (role-gated reverts); duplicate batch ids; missing evidence (an unattested certificate must not satisfy compliance); invalid transfers (recipient without role, self-transfer); the two-step handover (only the nominated recipient accepts); double attestation and self-attestation rejections; and revocation flipping compliance back while preserving the record.

Gas is measured directly from transaction receipts on the in-memory chain, which makes numbers deterministic per commit. View costs are \texttt{eth\_estimateGas} on an \texttt{eth\_call}---free for off-chain readers, but the price another contract would pay and a clean complexity signal. The design alternatives of Section~\ref{sec:gas} are measured with single-purpose bench contracts that keep everything except the choice under test identical to the main contract: same struct fields, same event shapes, same validation logic.

\section{Gas and Scalability Analysis}\label{sec:gas}

\begin{table}[!t]
\caption{Measured gas on the in-memory chain (solc 0.8.28, optimizer 200)}
\label{tab:gas}
\centering
\begin{tabular}{@{}l r@{}}
\toprule
Operation & Gas \\
\midrule
deploy contract & 3,187,558 \\
grantRole (per holder) & 51,495 \\
registerComponent (31-char id) & 119,946 \\
registerComponent (40-char id) & 164,851 \\
addCertificate & 221,392--224,204 \\
attestCertificate & 57,270 \\
revokeCertificate & 51,600 \\
assembleHBOM (6 components) & 352,534 \\
assembleHBOM (16 components) & 678,456 \\
initiateTransfer / acceptTransfer & 37,770 / 37,910 \\
verifyCompliance (6, view) & 205,538 \\
verifyCompliance (16, view) & 502,165 \\
getBoardLineage (6, view) & 546,811 \\
\midrule
seed world (30 parts, 89 certs, 1 board) & 28,918,618 \\
\bottomrule
\end{tabular}
\end{table}

Table~\ref{tab:gas} gives the operating profile; three findings drive the design discussion.

\subsection{Certificates Dominate}
Composing the seeded demo world costs 28.9M gas, and 69\% of it is \texttt{addCertificate}: 89 documents at roughly 224k each, each writing a full struct (two strings, nine value fields) plus an index append. The registry itself is cheap. \emph{Any} meaningful cost reduction must therefore target the evidence model, not the id keys or the custody machinery.

\subsection{Human-Readable Ids Are Nearly Free---Below 32 Characters}
Ids of 10, 20, and 31 characters all register within 252 gas of each other, because Solidity packs short strings into a single storage slot; the measured premium of a string key over a \texttt{bytes32} key in an otherwise identical bench contract is 4.6k gas ($+3.8\%$)---calldata hashing, not storage. The cliff is real but avoidable: at 40 characters the id spills into a second data slot and registration costs $+37\%$. The design response is a naming convention (all dataset ids are $\le$~31 characters), not a type change. Strings do cost on assembly---each \texttt{string[]} member is a pointer to separately stored data---making a six-component assembly $22\%$ more expensive than its \texttt{bytes32[]} equivalent.

\subsection{The Merkle Alternative: 87\% Cheaper, and Wrong for This Product}
A bench contract storing a single Merkle root per board assembles for 45,747 gas ($-87\%$) and verifies membership in 27,750 gas, nearly flat in HBOM size (27,750 $\to$ 29,000 from 6 to 16 leaves, versus a 205k $\to$ 502k linear scan). If the only question were ``is \emph{this exact} HBOM the one that was certified,'' this is the correct architecture. But a root cannot answer \emph{which} documents are missing, ends per-certificate history, and turns revocation into a tree rebuild. The buyer badge and the auditor queue are the product; we keep per-component records and accept the cost, noting a production hybrid (per-component roots, off-chain proofs, indexer-driven enumeration) in Section~\ref{sec:future}.

\subsection{AccessControl Is Worth Its Premium}
Against a flat \texttt{mapping(address => bool)} bench, OpenZeppelin enrollment costs $+7.6$k gas (the \texttt{RoleGranted} event and admin-branch check) and each gated call costs $+391$ gas ($+0.3\%$). For that premium the contract gains a role-admin hierarchy, introspection, and standardized governance events---the audit trail the system exists to produce. The scalability story is bounded and honest: \texttt{MAX\_COMPONENTS} keeps assembly under 680k and an on-chain compliance read under 505k, and both scale linearly in components, which caps the per-board cost by construction.

\section{Discussion}\label{sec:discussion}

\subsection{Blockchain versus a Centralized Registry}
A single-operator web registry would implement every function in this system for a fraction of the engineering cost, at zero gas, with faster queries. It would also have a failure mode this project exists to remove: every other tier must trust the operator, who can edit history, lock out a participant, or favor its own tier. The blockchain binding is not about technology preference but about \emph{mutual custody of the record}: no participant can amend a batch link, an attestation, or a custody event after the fact, and no participant needs to trust another's database. The measured price of that guarantee is Table~\ref{tab:gas}: roughly 6.1M gas (about 0.12 ETH at 20\,gwei) to onboard and certify a six-component board with full evidence, which is why the prototype runs on a local chain and any production deployment assumes L2 fees or batching rather than mainnet costs.

\subsection{Privacy}
Everything on-chain is public by design: addresses, batch ids, document hashes, CIDs, custody order. Document \emph{contents} stay off-chain, and a keccak256 commitment leaks nothing invertible---but the linkage structure itself is commercial information. A real deployment must decide whether a tier-2 vendor's customer list is acceptable on a public chain. The standard answers are organizational (a permissioned chain among consortium members), cryptographic (salted hashes to break linkability), or architectural (per-tier subchains with cross-chain board links). Our prototype chooses public visibility because the buyer-verifiability story depends on it, and we flag this as the first thing a production review would change.

\subsection{Interoperability and Regulation}
The three certificate types are an enum tied to EU/US regimes that change on their own schedules---RoHS exemption lists, REACH candidate-list updates, FCC part revisions. The contract stores \emph{that a document of a type exists and was attested}; it deliberately does not encode regulation text, so a regulatory change invalidates nothing on-chain---it changes what auditors will attest going forward and may trigger revocations. Two harder boundaries remain. First, the \emph{oracle problem}: the chain binds a document hash, not the truth of the document; binding ids to physical batches and documents to accredited issuers is institutional process outside any ledger~\cite{guin}. Second, immutability sits uneasily with rights-to-erasure regimes such as GDPR; addresses are pseudonymous but linkable, and revocation-by-flag (rather than deletion) is the compromise this design shares with most deployed provenance systems.

\subsection{Limitations and Threats}
The system proves document integrity and custody \emph{order}; it does not prove physical authenticity, laboratory correctness, or shipment contents~\cite{guin}. \texttt{ADMIN} is a trust root: whoever holds it can grant roles, and the prototype's role-admin policy is a deploy-time convention. IPFS pinning is a liveness promise, not a durability guarantee---unpinned documents degrade the system to hash-only verification. The contract is not upgradeable (deliberately; redeploys and reseeds are acceptable at course scale), and enumeration relies on event replay rather than an indexer. None of these are hidden: the demo script narrates them, and each maps to a future-work item.

\subsection{Future Improvements}\label{sec:future}
Four, in order of measured payoff: (1) \emph{batched evidence}---a root per certificate family per batch, cutting the dominant 69\% cost term while keeping per-batch gap enumeration on-chain; (2) an \emph{event indexer} replacing the client's full log replay, which removes the read-side scaling ceiling entirely; (3) \emph{document versioning} with supersede semantics, extending the current attest/revoke pair; and (4) \emph{account abstraction} so participants join with passkeys instead of raw keys, which would also let the admin role rotate without seed-phrase ceremonies.

\section{Conclusion}
A board's provenance is currently an act of trust distributed across companies with no shared record. We built the alternative at course scale: one role-gated contract where registration, evidence, attestation, assembly, and custody are public, append-only, and independently checkable, with documents anchored by content hashes and stored off-chain. The measurements justify the design's economics---certificates dominate cost, readable ids are free below 32 characters, exact-gap reporting costs almost nothing, and the role machinery from an audited library costs fractions of a percent---and the discussion bounds the claim honestly: this is a shared record of documents and custody, not a certificate of physical truth.

\begin{thebibliography}{1}

\bibitem{guin}
U. Guin, D. DiMase, and M. Tehranipoor, ``Counterfeit integrated circuits: Detection, avoidance, and the challenges ahead,'' \emph{J. Electron. Test.}, vol. 30, no. 1, pp. 9--23, Feb. 2014.

\bibitem{cui}
P. Cui, J. Dixon, U. Guin, and D. Dimase, ``A blockchain-based framework for supply chain provenance,'' \emph{IEEE Access}, vol. 7, pp. 157113--157125, 2019.

\bibitem{toyoda}
K. Toyoda, P. T. Mathiopoulos, I. Sasase, and T. Ohtsuki, ``A novel blockchain-based product ownership management system (POMS) for anti-counterfeits in the post supply chain,'' \emph{IEEE Access}, vol. 5, pp. 17465--17477, 2017.

\bibitem{oz}
OpenZeppelin, ``AccessControl,'' \emph{OpenZeppelin Contracts documentation}. [Online]. Available: https://docs.openzeppelin.com/contracts/5.x/access-control

\end{thebibliography}

\end{document}
